\chapter{Scientific integrity, generative AI and reproducibility}

\section{Declaration of generative AI usage}
\guided{The use of AI tools to rephrase text, generate code, create figures, suggest outlines or summarize resources must be explicitly declared. The student remains fully responsible for the accuracy of the content and compliance with citation rules.}

\begin{quote}
Example statement: \textit{Generative AI tools were used at certain stages to support text rephrasing, suggest a document structure or propose code examples. All content was subsequently verified, corrected and validated by the author, who assumes full scientific and academic responsibility for the final work.}
\end{quote}

\section{Minimum reproducibility expectations}
\begin{itemize}[leftmargin=1.2cm]
\item Datasets: source, size, date, acquisition mode, restrictions.
\item Tools: software versions, libraries, hardware environment, system configuration.
\item Critical parameters: hyperparameters, thresholds, frequencies, dimensions, protocols.
\item Validation elements: scripts, test notebooks, configuration captures, schematics, bills of materials.
\item Confidentiality: distinguish what is shareable, masked or anonymized.
\end{itemize}

\section{Compliance reminders}
\begin{itemize}[leftmargin=1.2cm]
\item \okitem{} Every borrowed figure, table, code excerpt or idea must be attributed.
\item \okitem{} Results must be interpreted critically.
\item \warnitem{} An excessively high similarity rate may compromise report acceptability.
\item \warnitem{} Fluent text is not scientific evidence if it is not linked to a method and verifiable results.
\end{itemize}
